\documentclass[a4paper,11pt]{article}

\usepackage[all]{carnotex}
\level{notes}

\title{Test complet CarnoTeX}
\author{}
\date{\today}

\begin{document}

\maketitle
\tableofcontents

\section{carnotex}

\cadre{CarnoDef}[Cadre avec titre]{
Contenu du cadre.
}

\cadre{CarnoProp}{
Cadre sans titre.
}

\Cbox{CarnoIdea}{Bloc documentaire}{
Contenu de la Cbox.
}

\colorline[CarnoProp]

\colorline*[CarnoThm][6cm][1pt]

\histo{
Texte historique.
}

\exobandeau{Exercice simple}{2}

Énoncé.

\exobandeau[5]{Exercice noté}{4}

Énoncé.

\col{
Colonne gauche.
}{
Colonne droite.
}

\section{carnomath}

\df{Définition non étoilée.}

\df*{Définition étoilée.}

\prop{Propriété non étoilée.}

\prop*{Propriété étoilée.}

\thrm{Théorème non étoilé.}

\thrm*{Théorème étoilé.}

\lem{Lemme non étoilé.}

\lem*{Lemme étoilé.}

\coro{Corollaire non étoilé.}

\coro*{Corollaire étoilé.}

\exe{Exemple.}

\exo{Exercice.}

\prb{Problème.}

\rem{Remarque.}

\dem{
Démonstration.

\cqfd
}

\sol{Solution.}

\[
\N,\quad \Z,\quad \D,\quad \Q,\quad \R,\quad \Cp
\]

\[
\vc{AB}
\qquad
\coord{2}{3}
\qquad
\coor{1}{2}{3}
\qquad
\vct{u}{2}{-1}
\]

\Oij

\Oijk

\section{carnopeda}

\level{notes}

\obj{
Objectifs.
}

\pre{
Prérequis.
}

\idee{
Idée pédagogique.
}

\piege{
Piège classique.
}

\retenir{
À retenir.
}

\note[Préparation]{
Note personnelle.
}

\source[Source]{
Source.
}

\experience[Retour]{
Retour d'expérience.
}

\todo[À faire]{
Tâche à faire.
}

\temps{10}

\section{Utilitaires}

\cqfd

\siecle{21}

\siecle*{3}

\rl

\end{document}